% Fragmento integrable. Preambulo anfitrion: amsmath, amssymb, longtable, hyperref.
\section{Classification and continuation of the first-root selector}
\label{hmtfg:fragmento:clasificacion-y-continuacion-del-selector-de-primeras-raices}

\subsection{Object and provenance}
\label{hmtfg:clasificacion:1}

This appendix preserves the family rendered by the HMT dodecaphase reader:

\[
u_0(x)=\frac1{12}\sum_{m=1}^{12}\left(1-\cos\frac{2(3m-1)x}{\sqrt{899}}\right),
\qquad f_\lambda(x)=1-\lambda u_0(x).
\]

The family originates in the operator and analytic antecedents presented in this development. The existence and compactness of its realizations of the infinite itinerary were proved in the subsections on realization of the word and preservation of its signs.

The following argument is original and uses only those properties and the return operations of the fixed family. It does not use a parameter from another family or the value of a universal constant as an input. The sheets of the itineraries used here are the two monotone branches of the critical reader; the remaining coordinates and HMT memory retain their residence in the state from which the reader originates.

We write
\[
c_j(\lambda)=f_\lambda^j(0),\qquad
\tau_j=(-1)^{v_2(j)},\qquad
W_n=\tau_1\cdots\tau_{2^n-1}.
\]
The itinerary of a critical orbit that returns to the origin for the first time ends in the zero symbol. The itinerary order is the unimodal order with a maximum: before comparing the symbol of index \(j\), the orientation factor is
\[
O_{j-1}(K)=\prod_{i=1}^{j-1}(-K_i).
\]
The ordinary order of the symbols is \(-1<0<1\).


\subsection{Classification before the infinite word}
\label{hmtfg:clasificacion:2}

\textbf{Classification theorem.} Let \(F:I\to I\), with \(I=[\ell,1]\), \(\ell<0\), be a continuous map, strictly increasing to the left of zero and strictly decreasing to its right, whose maximum satisfies \(F(0)=1\). Suppose zero has exact period \(2^n\), with \(n\ge1\), and its critical itinerary \(K\) satisfies \(K<\tau\). Then
\[
\boxed{K=W_n0.}
\]
The conclusion concerns the itinerary and allows distinct parameters in a family to have that same itinerary.

\textbf{Proof.} Write \(d_j=F^j(0)\). For period two the itinerary is \(+0=W_10\). Consider an exact period \(p=2^n\ge4\).

We have \(d_2<0\). Indeed, \(d_2=0\) would give period two. If \(d_2>0\), monotonicity of the positive branch gives
\[
F([d_2,1])=[d_2,d_3]\subset[d_2,1],
\]
and the critical orbit remains in a strictly positive interval, incompatible with a return to zero.

The initial signs are then \(+,-\). The orientation factor for comparing the third sign is \(-1\); hence \(K<\tau\) forces \(d_3>0\). A zero at that position would produce period three.

Nor can \(d_4<0\) occur. Under this hypothesis,
\[
d_2<d_4<0,
\qquad F([d_2,d_4])=[d_3,d_5]\subset(0,1],
\]
since \(d_4=F(d_3)>F(1)=d_2\) and \(d_5=F(d_4)>F(d_2)=d_3\). The positive branch is decreasing, so that
\[
F^2([d_2,d_4])=[d_6,d_4]\subset[d_2,d_4].
\]
Here \(d_6\ge d_2\), because \(d_5\le1\), and \(d_6<d_4\), because \(d_5>d_3\). This strictly negative interval traps the even returns of the critical orbit and excludes a return to zero. Therefore \(d_4\ge0\). If \(p=4\), one obtains precisely \(+-+0=W_20\).

Now suppose \(p\ge8\). Then \(d_4>0\). The orientation factors for the fifth and sixth signs are, respectively, \(-1\) and \(+1\). Since \(\tau_5=+1\) and \(\tau_6=-1\), the inequality \(K<\tau\) forces
\begin{equation}
\operatorname{sign}(d_1,\ldots,d_6)=(+,-,+,+,+,-).
\label{hmtfg:clasificacion:eq:1}
\end{equation}
Zero signs at those positions are excluded by the exact period.

We construct the return interval
\[
J=[d_2,d_4].
\]
The inequality \(F(d_3)=d_4>0>d_6=F(d_5)\), together with \(d_3,d_5>0\), implies \(d_3<d_5\); therefore
\[
F(J)=[d_3,1].
\]
Moreover, \(d_4<d_3\). To check this, the image under \(F\) of the interval with endpoints \(d_3,d_4\) is strictly positive, since its endpoint images are \(d_4,d_5>0\). On that interval \(F^2\) is strictly increasing. Its endpoint values satisfy
\[
F^2(d_3)=d_5>0>d_6=F^2(d_4),
\]
which forces \(d_3>d_4\). Consequently,
\begin{equation}
J\cap F(J)=\varnothing,
\qquad F^2(J)=F([d_3,1])=[d_2,d_4]=J.
\label{hmtfg:clasificacion:eq:2}
\end{equation}

The map \(F^2|_J\) has a unique minimum at zero. Conjugacy by \(h(x)=d_2x\), which reverses orientation, defines
\begin{equation}
G=h^{-1}\circ F^2\circ h:
\left[\frac{d_4}{d_2},1\right]\longrightarrow
\left[\frac{d_4}{d_2},1\right].
\label{hmtfg:clasificacion:eq:3}
\end{equation}
This map satisfies the same maximum hypotheses as \(F\), with \(G(0)=1\), and its critical point has exact period \(p/2\). Its itinerary is
\begin{equation}
K'_j=-K_{2j}.
\label{hmtfg:clasificacion:eq:4}
\end{equation}
All odd-indexed symbols of \(K\) are positive, because the even positions belong to \(J\) and the odd positions to \(F(J)\subset(0,1]\).

The transformation \eqref{hmtfg:clasificacion:eq:4} preserves the comparison order with \(\tau\). Indeed, \(\tau_{2j-1}=+1\) and \(\tau_{2j}=-\tau_j\). If \(q\) is the first position at which \(K'\) and \(\tau\) differ, the first position of difference between \(K\) and \(\tau\) is \(2q\), and
\begin{equation}
O_{2q-1}(K)=-O_{q-1}(K'),
\qquad K_{2q}-\tau_{2q}=-(K'_q-\tau_q).
\label{hmtfg:clasificacion:eq:5}
\end{equation}
The product determining the order is identical. The identities remain valid when the first distinct symbol is the final zero. Thus \(K'<\tau\).

Induction applied to \(G\) gives \(K'=W_{n-1}0\). The positive odd positions and \eqref{hmtfg:clasificacion:eq:4} reconstruct exactly \(K=W_n0\). The theorem is proved. \(\square\)

The theorem also proves that each classified map has the restrictive doubling returns preceding its superstable return, with the domains of \eqref{hmtfg:clasificacion:eq:2}–\eqref{hmtfg:clasificacion:eq:3}. It does not presuppose monotonicity of a family parameter.

\textbf{Symmetric normalization of the fixed family.} When \(F\) is even and defined on \([-1,1]\), the preceding inequalities additionally give
\[
F(d_4)=d_5>d_3=F(d_2)=F(-d_2).
\]
Since both arguments \(d_4,-d_2\) are positive and \(F\) decreases on that branch,
\begin{equation}
0<d_4<-d_2<1.
\label{hmtfg:clasificacion:eq:5a}
\end{equation}
Therefore \(F^2([d_2,-d_2])=[d_2,d_4]\), and the same conjugacy \(h(x)=d_2x\) defines on all of \([-1,1]\) an even map with a maximum whose image is \([d_4/d_2,1]\subset[-1,1]\). It agrees exactly with the operator \(RF=-F(F(-ax))/a\), \(a=-F(1)\), used in the corpus. This symmetric extension preserves the critical return, its period and its itinerary.


\subsection{First parameter with infinite itinerary}
\label{hmtfg:clasificacion:3}

We work on the compact interval already used in the existence proof,
\[
L=\left[\frac1{2u_0(1)},\frac2{u_0(1)}\right].
\]
Let
\[
\Lambda_\tau=\{\lambda\in L:
\operatorname{sign}c_j(\lambda)=\tau_j\text{ for every }j\ge1\}.
\]
By the existence and sign-separation proofs of the prolongation, this set is nonempty and compact. Therefore there exists
\begin{equation}
a=\min\Lambda_\tau.
\label{hmtfg:clasificacion:eq:6}
\end{equation}

\textbf{Comparison lemma.} For every \(\lambda\in[1/(2u_0(1)),a)\), we have \(K_\lambda<\tau\).

\textbf{Proof.} At the left endpoint, the itinerary is \(+^\infty<\tau\). The separation argument for realization of the infinite word proves that the strict comparison sets with \(\tau\) are open: in an infinite itinerary distinct from \(\tau\), the first finite difference persists by continuity; at a superstable parameter, the two neighboring periodic words and the word with a final zero lie on the same side of \(\tau\), by its strict maximality. By definition, the interval preceding \(a\) contains no parameters with itinerary \(\tau\). Its connectedness maintains the initial comparison throughout that interval. \(\square\)


\subsection{Existence of each first new root}
\label{hmtfg:clasificacion:4}

We retain the original selector. The first root is
\[
\lambda_1=\frac1{u_0(1)},
\]
and, for \(n\ge2\), \(\lambda_n\) is the first root of \(c_{2^n}\) to the right of \(\lambda_{n-1}\) satisfying
\begin{equation}
c_{2^j}(\lambda_n)\ne0\quad(1\le j<n).
\label{hmtfg:clasificacion:eq:7}
\end{equation}
Since every first period dividing \(2^n\) is a power of two, \eqref{hmtfg:clasificacion:eq:7} is equivalent to requiring exact critical period \(2^n\).

\textbf{Existence lemma.} All roots of this selector exist and satisfy
\begin{equation}
\lambda_1<\lambda_2<\cdots<\lambda_n<a.
\label{hmtfg:clasificacion:eq:8}
\end{equation}

\textbf{Proof.} The sign \(c_2(a)<0\) implies \(\lambda_1<a\). Suppose \(\lambda_{n-1}<a\) has been constructed, and put \(p=2^{n-1}\). The analytic function \(c_p(\lambda)\) is not identically zero, since \(c_p(a)\ne0\). Its zero set in \([\lambda_{n-1},a]\) is finite, nonempty and does not contain \(a\). Let \(z<a\) be its last zero.

On \((z,a]\), the sign of \(c_p\) is constant and equal to \(s=\tau_p\). For \(g_\lambda=f_\lambda^p\), we have
\[
g_\lambda(0)=c_p(\lambda),\qquad g'_\lambda(0)=0.
\]
A uniform local bound on the second derivative gives
\begin{equation}
c_{2p}(\lambda)
=g_\lambda(c_p(\lambda))
=c_p(\lambda)+O(c_p(\lambda)^2).
\label{hmtfg:clasificacion:eq:9}
\end{equation}
As \(\lambda\) approaches \(z\) from the right, the first-order term dominates; hence \(c_{2p}(\lambda)\) has sign \(s\). At \(a\), its sign is
\[
\tau_{2p}=-\tau_p=-s.
\]
The intermediate value theorem provides a zero \(\beta\in(z,a)\) of \(c_{2p}\). Since \(c_p\ne0\) on that interval, the critical period at \(\beta\) is exactly \(2p\).

Analyticity and \(c_{2p}(a)\ne0\) imply that the zeros of \(c_{2p}\) in \([\lambda_{n-1},a]\) are finite. Among its zeros of exact period \(2p\) to the right of \(\lambda_{n-1}\), there is therefore a first one. It is the original selector \(\lambda_n\), and it satisfies \(\lambda_n\le\beta<a\). \(\square\)

Use of the last auxiliary zero \(z\) belongs solely to the existence proof. The selected parameter remains the \textbf{first new root} \(\lambda_n\).


\subsection{Convergence of the global selector to the first infinite itinerary}
\label{hmtfg:clasificacion:5}

\textbf{Global continuation theorem.} For the fixed HMT family and the original selector, all selected itineraries are canonical and
\begin{equation}
\boxed{K_{\lambda_n}=W_n0,\qquad
\lambda_n\nearrow\min\Lambda_\tau.}
\label{hmtfg:clasificacion:eq:10}
\end{equation}

\textbf{Proof.} The comparison and existence lemmas place each selected parameter below \(a\), with an itinerary less than \(\tau\). The classification theorem identifies that itinerary with \(W_n0\). The sequence is increasing and bounded by \(a\); let \(b\le a\) be its limit.

Fix \(p\ge1\). For \(n\) sufficiently large, \(2p<2^n\). Then the signs of \(c_p(\lambda_n)\) and \(c_{2p}(\lambda_n)\) agree with \(\tau_p\) and \(\tau_{2p}=-\tau_p\), respectively. The uniform bound on second derivatives,
\[
\|(f_\lambda^p)''\|\le B_p,
\qquad B_p=16^p\frac{16^p-1}{15},
\]
and the identity \((f_\lambda^p)'(0)=0\) imply
\[
|c_{2p}(\lambda_n)-c_p(\lambda_n)|
\le\tfrac12B_p|c_p(\lambda_n)|^2.
\]
Since the two terms on the left-hand side have opposite signs,
\begin{equation}
|c_p(\lambda_n)|>2/B_p.
\label{hmtfg:clasificacion:eq:11}
\end{equation}
By continuity, \(c_p(b)\) retains the sign \(\tau_p\) and satisfies \(|c_p(b)|\ge2/B_p>0\). This holds for each \(p\), so that \(b\in\Lambda_\tau\). Minimality of \(a\) gives \(b\ge a\), and therefore \(b=a\). \(\square\)

The complete proof preserves the global selector and produces its itineraries and its limit. The information from the first eight levels is compatible with the result, but the proof of \eqref{hmtfg:clasificacion:eq:10} does not extrapolate those eight levels: it uses induction, invariant intervals and uniform separation of every sign.


\subsection{Logical checks and falsifiers}
\label{hmtfg:clasificacion:7}

\par\noindent\textbf{1.}\enspace The classification theorem requires a strict unimodal maximum and an invariant interval. If either branch loses monotonicity, the intervals \eqref{hmtfg:clasificacion:eq:2} may cease to be restrictive and the classification must be checked again.
\par\noindent\textbf{2.}\enspace The condition \(K<\tau\) is essential. Beyond the limit word, centers of period \(2^n\) with other combinatorics may exist; the theorem does not identify them with \(W_n0\).
\par\noindent\textbf{3.}\enspace Analyticity in the parameter provides isolation and finiteness of the zeros on the compact set. For a merely continuous family, the existence of a strictly subsequent first root requires an additional argument.
\par\noindent\textbf{4.}\enspace A limit of critical roots may, in general, lose signs upon reaching a return of lower period. Inequality \eqref{hmtfg:clasificacion:eq:11} excludes precisely that phenomenon for this sequence and each fixed horizon.
\par\noindent\textbf{5.}\enspace Convergence of the first roots to the first parameter with itinerary \(\tau\) is not equivalent to uniqueness of all parameters with that itinerary and does not by itself determine a metric rate. Transverse crossing and branch identification provide those subsequent conclusions.
